% Auto-generated by analysis/export_edges.py. Do not edit by hand.
\begin{table*}[t]
\caption{The complete edge set of $K_G$. Every edge asserts that acting on the source is a guideline-recognised route to changing the target. Grade is the value the ADA Professional Practice Committee assigned to the cited recommendation; \emph{narrative} marks an edge supported by section text rather than a numbered recommendation, and is not graded.}
\label{tab:edges-full}
\centering\scriptsize
\begin{tabular}{@{}lllp{9.2cm}@{}}
\toprule
\textbf{Edge} & \textbf{Grade} & \textbf{ADA source} & \textbf{What the recommendation licenses} \\
\midrule
\texttt{BMI} $\rightarrow$ \texttt{HighBP} & A & S8, Rec 8.5 & Weight loss of 5-7\% improves glycaemia and intermediate cardiovascular risk factors. \\
\texttt{BMI} $\rightarrow$ \texttt{HighBP} & A & S10, Rec 10.5 & Weight loss is the first lifestyle behaviour advised above 120/80 mmHg. \\
\texttt{BMI} $\rightarrow$ \texttt{HighChol} & A & S8, Rec 8.5 & Same recommendation; dyslipidaemia is among the intermediate cardiovascular risk factors. \\
\texttt{Fruits} $\rightarrow$ \texttt{BMI} & A & S5, Rec 5.12 & Nutrition is one of the three components of the weight-management plan. \\
\texttt{Fruits} $\rightarrow$ \texttt{HighBP} & A & S10, Rec 10.5 & A DASH-style eating pattern is one of the advised lifestyle behaviours. \\
\texttt{HvyAlcoholConsump} $\rightarrow$ \texttt{HighBP} & A & S10, Rec 10.5 & Limiting or avoiding alcohol is one of the advised lifestyle behaviours. \\
\texttt{PhysActivity} $\rightarrow$ \texttt{BMI} & A & S5, Rec 5.12 & A weight-management plan built on nutrition, physical activity and behavioural support target... \\
\texttt{PhysActivity} $\rightarrow$ \texttt{BMI} & A & S8, Rec 8.8a & Structured counselling centred on nutrition and physical activity to achieve a daily energy d... \\
\texttt{PhysActivity} $\rightarrow$ \texttt{HighBP} & A & S10, Rec 10.5 & Increased physical activity is one of the advised lifestyle behaviours. \\
\texttt{Smoker} $\rightarrow$ \texttt{HighBP} & A & S10, Rec 10.5 & Smoking cessation is one of the advised lifestyle behaviours. \\
\texttt{Veggies} $\rightarrow$ \texttt{BMI} & A & S5, Rec 5.12 & Nutrition is one of the three components of the weight-management plan. \\
\texttt{Veggies} $\rightarrow$ \texttt{HighBP} & A & S10, Rec 10.5 & A DASH-style eating pattern is one of the advised lifestyle behaviours. \\
\texttt{Fruits} $\rightarrow$ \texttt{BMI} & B & S5, Rec 5.14 & Eating patterns should emphasise whole fruits, among other principles. \\
\texttt{HvyAlcoholConsump} $\rightarrow$ \texttt{HighBP} & B & S5, Rec 5.18 & Adults with or at risk of diabetes should not exceed recommended daily alcohol limits. \\
\texttt{Veggies} $\rightarrow$ \texttt{BMI} & B & S5, Rec 5.14 & Eating patterns should emphasise nonstarchy vegetables, among other principles. \\
\texttt{Fruits} $\rightarrow$ \texttt{HighChol} & C & S10, Rec 10.15 & Lifestyle therapy is intensified for raised triglycerides or low HDL. \\
\texttt{PhysActivity} $\rightarrow$ \texttt{HighChol} & C & S10, Rec 10.15 & Lifestyle therapy is intensified for raised triglycerides or low HDL. \\
\texttt{Veggies} $\rightarrow$ \texttt{HighChol} & C & S10, Rec 10.15 & Lifestyle therapy is intensified for raised triglycerides or low HDL. \\
\texttt{BMI} $\rightarrow$ \texttt{GenHlth} & narrative & S8, obesity narr. & Obesity treatment is framed as improving health outcomes broadly. \\
\texttt{BMI} $\rightarrow$ \texttt{PhysHlth} & narrative & S8, obesity narr. & Obesity treatment is framed as improving health outcomes broadly. \\
\texttt{HighBP} $\rightarrow$ \texttt{GenHlth} & narrative & S10, narr. & Blood-pressure control is framed as reducing cardiovascular risk and improving outcomes. \\
\texttt{HighChol} $\rightarrow$ \texttt{GenHlth} & narrative & S10, narr. & Lipid control is framed as reducing cardiovascular risk and improving outcomes. \\
\texttt{PhysActivity} $\rightarrow$ \texttt{GenHlth} & narrative & S5, physical activity narr. & Physical activity is described as favourably affecting glycaemia, cardiovascular risk factors... \\
\texttt{PhysActivity} $\rightarrow$ \texttt{MentHlth} & narrative & S5, physical activity narr. & Physical activity is described as favourably affecting psychosocial well-being, and as benefi... \\
\texttt{PhysActivity} $\rightarrow$ \texttt{PhysHlth} & narrative & S5, physical activity narr. & Physical activity is described as favourably affecting mobility and physical function. \\
\texttt{Smoker} $\rightarrow$ \texttt{PhysHlth} & narrative & S5, tobacco narr. & Smoking is described as causally linked to multiple health risks that affect morbidity in peo... \\
\bottomrule
\end{tabular}
\end{table*}
